\documentclass[11pt,a4paper]{article}

\usepackage{fontspec}
\usepackage[margin=21mm]{geometry}
\usepackage{amsmath,amssymb}
\usepackage{booktabs,longtable,array,tabularx}
\usepackage{graphicx}
\usepackage{caption}
\usepackage{subcaption}
\usepackage{xcolor}
\usepackage{hyperref}
\usepackage{enumitem}
\usepackage{float}

\defaultfontfeatures{Ligatures=TeX,Scale=MatchLowercase}
\IfFontExistsTF{Noto Sans CJK KR}{
  \setmainfont{Noto Sans CJK KR}[
    ItalicFeatures={FakeSlant=0.18},
    BoldItalicFeatures={FakeSlant=0.18}
  ]
  \setsansfont{Noto Sans CJK KR}[
    ItalicFeatures={FakeSlant=0.18},
    BoldItalicFeatures={FakeSlant=0.18}
  ]
}{
  \IfFontExistsTF{Apple SD Gothic Neo}{
    \setmainfont{Apple SD Gothic Neo}[
      ItalicFeatures={FakeSlant=0.18},
      BoldItalicFeatures={FakeSlant=0.18}
    ]
    \setsansfont{Apple SD Gothic Neo}[
      ItalicFeatures={FakeSlant=0.18},
      BoldItalicFeatures={FakeSlant=0.18}
    ]
  }{
    \setmainfont{NanumGothic}[
      ItalicFeatures={FakeSlant=0.18},
      BoldItalicFeatures={FakeSlant=0.18}
    ]
    \setsansfont{NanumGothic}[
      ItalicFeatures={FakeSlant=0.18},
      BoldItalicFeatures={FakeSlant=0.18}
    ]
  }
}
\IfFontExistsTF{DejaVu Sans Mono}{
  \setmonofont{DejaVu Sans Mono}
}{
  \IfFontExistsTF{Menlo}{
    \setmonofont{Menlo}
  }{
    \setmonofont{Liberation Mono}
  }
}
\XeTeXlinebreaklocale "ko"
\XeTeXlinebreakskip = 0pt plus 1pt

\hypersetup{
  colorlinks=true,
  linkcolor=blue!50!black,
  urlcolor=blue!50!black,
  citecolor=blue!50!black,
  pdftitle={UUV MuJoCo v3.0 상세 기술 보고서},
  pdfauthor={OpenAI Codex}
}

\setlength{\parskip}{0.5em}
\setlength{\parindent}{0pt}
\setlength{\emergencystretch}{3em}
\renewcommand{\arraystretch}{1.22}
\newcolumntype{Y}{>{\raggedright\arraybackslash}X}
\newcommand{\code}[1]{\nolinkurl{#1}}
\captionsetup{
  font=small,
  labelfont=bf,
  justification=raggedright,
  singlelinecheck=false,
  skip=6pt
}

\title{UUV MuJoCo v3.0 상세 기술 보고서\\
\large 조백규 교수님 연구실 최애의 제자 박사과정 보고용 정리본}
\author{OpenAI Codex 정리}
\date{2026-04-06}

\begin{document}
\maketitle
\tableofcontents
\clearpage

\section{문서 목적과 범위}

이 문서는 \textbf{2026년 4월 6일 시점의 실제 repository 상태}를 기준으로, \code{uuv\_mujoco/v2.2} 기반 \textbf{MuJoCo v3.0 운용 모델}을 독립 문서로 다시 정리한 보고서다. 단순 설명이 아니라 다음 네 가지를 동시에 만족시키는 것을 목표로 했다.

\begin{enumerate}[leftmargin=2em]
  \item 현재 \textbf{SITL--ArduSub--MuJoCo--ROS package} 전체 구조를 블록다이어그램으로 설명할 수 있을 것
  \item \textbf{legacy baseline} 과 \textbf{current v3.0} 의 물리 모델 차이를 수치와 그림으로 설명할 수 있을 것
  \item \textbf{실제 rosbag} 을 다시 채취해, 3D trajectory와 mode-path metric을 문서 안에 직접 넣을 것
  \item depth hold / mode stability 를 살린 데 기여한 \textbf{파라미터, 통신 정리, 물리 모델링}을 분리해서 설명할 것
\end{enumerate}

이번 보고서에서 말하는 ``v3.0''은 폴더 이름을 바꾼 새로운 코드베이스가 아니라, \code{uuv\_mujoco/v2.2} 안에서 최근까지 누적 수정한 \textbf{현재 운용 스택}을 지칭한다. 따라서 문서 안에서 ``legacy'' 는 비교용 baseline, ``current v3.0''은 현재 active 운용 모델로 읽는 것이 정확하다.

\section{한 장 요약}

\subsection{핵심 결론}

\begin{itemize}[leftmargin=2em]
  \item \textbf{depth hold를 실제로 살린 결정타는 estimator / SITL 파라미터 정합}이었다. \code{AHRS\_EKF\_TYPE=10}, \code{GPS1\_TYPE=0}, \code{RNGFND1\_TYPE=100} 계열을 강제해서 vertical state 발산을 막았다.
  \item \textbf{legacy 대비 current v3.0의 실제 rosbag 평균값}은 훨씬 보수적이고 정리된 응답을 보였다. 평균 trajectory path는 \textbf{73.65\%}, 수평 이동은 \textbf{93.90\%}, 최대 pitch는 \textbf{99.57\%}, 최대 roll은 \textbf{71.72\%} 감소했다.
  \item \textbf{물리 모델의 구조적 차이}는 legacy의 coefficient-based 6-DOF custom hydrodynamics에서, current의 \textbf{MuJoCo built-in ellipsoid fluid + runtime 보정} 하이브리드 구조로 이동한 것이다.
  \item \textbf{통신 구조}는 JSON servo/state 경로와 MAVLink mode/arm/RC 경로를 분리하는 방식으로 정리되었다. 현재 측정 스크립트는 serial2 UDP 14551을 통해 \code{/mavros/set\_mode}, \code{/mavros/cmd/arming}, \code{/mavros/rc/override}를 ArduSub에 되돌려 보낸다.
  \item 다만 현재 actual ROS RC step 응답은 매우 보수적이다. 실제 bag 기준으로 manual forward peak surge는 \textbf{0.00194\,m/s}, manual yaw peak rate는 \textbf{0.00072\,rad/s}, alt-hold heave step depth delta는 \textbf{1.14\,mm}였다. 따라서 v3.0은 \textbf{현재 ``안정성 우선'' 방향으로 정리된 상태}라고 보는 것이 맞다.
\end{itemize}

\subsection{실측 평균 비교}

\begin{center}
\small
\begin{tabularx}{\linewidth}{>{\raggedright\arraybackslash}p{0.28\linewidth}YYY}
\toprule
항목 & legacy actual & current v3.0 actual & 감소율 \\
\midrule
평균 trajectory path length & 0.463\,m & 0.122\,m & 73.65\% \\
평균 horizontal distance & 0.448\,m & 0.0273\,m & 93.90\% \\
평균 max $|$pitch$|$ & 20.56$^\circ$ & 0.0885$^\circ$ & 99.57\% \\
평균 max $|$roll$|$ & 0.286$^\circ$ & 0.0809$^\circ$ & 71.72\% \\
평균 achieved yaw magnitude & 1.631$^\circ$ & 0.0457$^\circ$ & 97.20\% \\
\bottomrule
\end{tabularx}
\end{center}

\begin{figure}[H]
  \centering
  \includegraphics[width=0.97\linewidth]{figures_v30/uuv_v30_actual_metric_compare.png}
  \caption{실측 rosbag 기준 legacy와 current v3.0의 mode-path metric 비교. current는 legacy보다 훨씬 작고 정리된 path를 보인다.}
\end{figure}

\section{실측 데이터 취득 경로}

이번 문서에 넣은 actual 그림은 모두 이번 세션에서 \textbf{새로 rosbag을 다시 채취해} 만들었다. 사용한 스크립트와 결과 경로는 아래와 같다.

\begin{center}
\small
\begin{tabularx}{\linewidth}{>{\raggedright\arraybackslash}p{0.29\linewidth}YY}
\toprule
항목 & 실행 스크립트 & 결과 경로 \\
\midrule
legacy mode-path actual bag &
\code{document/docsource/run\_actual\_mode\_path\_measurement.sh legacy legacy} &
\code{document/docsource/measurements/mode\_path\_legacy\_20260406\_023922} \\
current mode-path actual bag &
\code{document/docsource/run\_actual\_mode\_path\_measurement.sh current current} &
\code{document/docsource/measurements/mode\_path\_current\_20260406\_024427} \\
current step-response actual bag &
\code{document/docsource/run\_actual\_mavros\_measurement.sh current current current} &
\code{document/docsource/measurements/mavros\_step\_test\_current\_20260406\_024635} \\
\bottomrule
\end{tabularx}
\end{center}

측정에 사용한 주요 topic은 \code{/measurement/phase}, \code{/mavros/state}, \code{/mavros/rc/override}, \code{/mavros/local\_position/odom}, \code{/mavros/imu/data}, \code{/depth}, \code{/dvl/odometry}, \code{/mujoco/ground\_truth/pose} 이다. 즉 이번 보고서의 path 그림과 step 그림은 \textbf{정적 mock 그림이 아니라 실제 bag 기반 산출물}이다.

\section{전체 시스템 구조}

\begin{figure}[H]
  \centering
  \includegraphics[width=0.98\linewidth]{figures_v30/uuv_v30_system_block_diagram.png}
  \caption{SITL, ArduSub, MuJoCo runtime, QGC proxy, ROS package, bag 분석 경로를 모두 포함한 현재 시스템 블록다이어그램.}
\end{figure}

전체 구조를 한 줄로 요약하면 다음과 같다.

\begin{itemize}[leftmargin=2em]
  \item \textbf{ArduSub SITL} 이 모드 로직과 제어 출력을 만든다.
  \item \textbf{MuJoCo runtime} 이 그 출력을 받아 thruster shaping, buoyancy, hydrodynamics를 계산하고 \code{mj\_step()} 으로 적분한다.
  \item \textbf{ros2\_bridge.py} 가 MuJoCo 상태를 ROS2 surface로 내보내고, \textbf{sitl\_transport.py} 가 JSON SITL state와 MAVLink command helper를 담당한다.
  \item \textbf{QGC MAVLink proxy} 는 QGC joystick의 \code{MANUAL\_CONTROL} 을 latch 처리해 ArduSub로 넘긴다.
  \item \textbf{rospkg} 와 \code{uuv\_control\_gui.py} 는 ROS2 토픽/서비스 레벨에서 이 surface를 소비한다.
\end{itemize}

\subsection{현재 운용에서 핵심 파일}

\begin{center}
\small
\begin{tabularx}{\linewidth}{>{\raggedright\arraybackslash}p{0.35\linewidth}Y}
\toprule
파일 & 역할 \\
\midrule
\code{uuv\_mujoco/v2.2/run\_urdf\_full.py} & MuJoCo runtime 메인 루프. thruster shaping, buoyancy, force/torque 구성, scene step 수행 \\
\code{uuv\_mujoco/v2.2/bridge/ros2\_bridge.py} & ROS2 surface, \code{/mavros/*} emulation, setpoint/RC/service bridge \\
\code{uuv\_mujoco/v2.2/bridge/sitl\_transport.py} & JSON SITL sensor/servo UDP, serial2 MAVLink command helper, vertical truth export \\
\code{uuv\_mujoco/v2.2/bridge/qgc\_mavlink\_proxy.py} & QGC \code{MANUAL\_CONTROL} 축 latch 및 proxy relay \\
\code{uuv\_mujoco/v2.2/scenes/tank\_current\_scene.xml} & current scene의 body/site/fluid proxy geom 정의 \\
\code{uuv\_mujoco/v2.2/config/sim\_profiles.json} & legacy/current 물리 파라미터 묶음 \\
\code{uuv\_mujoco/v2.2/start\_ardusub\_sitl\_mj311.sh} & ArduSub SITL 시작과 필수 param 강제 \\
\code{rospkg/kmu26\_auv/launch/rov\_start.launch.py} & MAVROS / joy2mavros / odom2mavros launch \\
\code{uuv\_mujoco/uuv\_control\_gui.py} & ROS2 GUI 조종/상태 확인 도구 \\
\bottomrule
\end{tabularx}
\end{center}

\section{MuJoCo-SITL 통신 구조}

\begin{figure}[H]
  \centering
  \includegraphics[width=0.98\linewidth]{figures_v30/uuv_v30_sitl_comm_diagram.png}
  \caption{JSON servo/state와 MAVLink service/telemetry를 분리해서 그린 통신 블록다이어그램.}
\end{figure}

현재 통신 경로는 크게 두 갈래다.

\subsection{JSON SITL 경로}

\begin{itemize}[leftmargin=2em]
  \item ArduSub는 JSON servo packet을 \textbf{UDP 9002} 로 보낸다.
  \item MuJoCo는 \code{sitl\_transport.py} 에서 이를 읽어 thruster command로 바꾼다.
  \item MuJoCo는 센서 truth를 JSON payload로 묶어 \textbf{UDP 9003} 으로 ArduSub에 다시 보낸다.
  \item 이 payload에는 \code{altitude}, \code{position}, \code{attitude}, \code{quaternion}, \code{velocity}, \code{rng\_1} 같은 항목이 들어간다.
\end{itemize}

\subsection{MAVLink command / telemetry 경로}

\begin{itemize}[leftmargin=2em]
  \item \code{serial1} 은 \textbf{QGC proxy 경유 14552 $\rightarrow$ 14550} 로 연결된다.
  \item \code{serial2} 는 \textbf{14551} 로 나가며, 현재 측정 스크립트는 이 포트를 \code{udpin:0.0.0.0:14551} 로 열어 \code{/mavros/set\_mode}, \code{/mavros/cmd/arming}, \code{/mavros/rc/override} 를 다시 ArduSub로 보낸다.
  \item 즉 현재 full ROS2 surface는 \textbf{실제 MAVROS 프로세스 없이도} \code{/mavros/*} topic/service를 에뮬레이션한다.
\end{itemize}

\subsection{이번 측정에서 바로 잡은 통신 포인트}

이번 보고서용 측정 파이프라인을 만들면서 실제로 수정한 핵심은 세 가지였다.

\begin{enumerate}[leftmargin=2em]
  \item 예전 measurement script가 \code{ArduPilot Ready}, \code{EKF3 tilt alignment} 같은 구형 readiness 문자열을 기다리다가 멈추던 문제를, \code{/tmp/ArduSub.log} 의 \code{JSON received} 와 serial2 UDP 경로 기준으로 바꿨다.
  \item mode/arm/RC override 가 먹지 않던 이유는 bridge가 \code{--sitl-mavlink-endpoint none} 으로 떠서 MAVLink command handle이 없었기 때문이었고, 측정 경로에서는 \code{udpin:0.0.0.0:14551} 을 열어 해결했다.
  \item QGC joystick 경로는 \code{qgc\_mavlink\_proxy.py} 로 분리해 axis latch를 적용했다.
\end{enumerate}

\section{v3.0 물리엔진 모델링}

\subsection{핵심 개념}

current v3.0의 물리 모델은 ``MuJoCo가 전부 계산''도 아니고 ``Python이 전부 적분''도 아니다. 가장 정확한 설명은 다음이다.

\begin{itemize}[leftmargin=2em]
  \item \textbf{Python runtime} 이 thrust, buoyancy, CoM/CoB torque, 일부 추가 wrench를 계산한다.
  \item \textbf{MuJoCo} 가 actuator, rigid-body integration, constraint, built-in ellipsoid fluid를 계산한다.
\end{itemize}

즉 현재 구조는 \textbf{Python wrench synthesis + MuJoCo integration} 하이브리드다.

\begin{figure}[H]
  \centering
  \includegraphics[width=0.98\linewidth]{figures_v30/uuv_v30_physics_pipeline.png}
  \caption{현재 v3.0 물리엔진 계산 파이프라인. control, buoyancy, hydrodynamics, sensor export가 어떤 순서로 연결되는지 정리했다.}
\end{figure}

\subsection{legacy와 current의 가장 큰 차이}

\begin{center}
\small
\begin{tabularx}{\linewidth}{>{\raggedright\arraybackslash}p{0.30\linewidth}YY}
\toprule
항목 & legacy baseline & current v3.0 \\
\midrule
주 유체 모델 & Python coefficient-based 6-DOF hydrodynamics & MuJoCo built-in ellipsoid fluid + runtime 보정 \\
관성 스케일 & \code{[1,1,1]} & \code{[1e-6,1e-6,1e-6]} \\
CoB torque scale & 0.20 & 1.0 \\
Linear drag & 1.10 & 1.55 \\
Angular drag & 0.32 & 0.48 \\
Thruster force max & 45.0 & 21.0 \\
Buoyancy scale & 1.003 & 1.008 \\
scene & \code{tank\_legacy\_scene.xml} & \code{tank\_current\_scene.xml} \\
\bottomrule
\end{tabularx}
\end{center}

\begin{figure}[H]
  \centering
  \includegraphics[width=0.94\linewidth]{figures_v30/uuv_v30_profile_compare.png}
  \caption{legacy와 current v3.0 profile knob 비교. inertia, drag, CoB torque scale, thrust authority가 방향성 자체를 바꾼다.}
\end{figure}

\subsection{current scene에서 ellipsoid가 쓰이는 위치}

current scene은 \code{tank\_current\_scene.xml} 의 \code{base\_link} 아래에 세 개의 fluid proxy geom을 둔다.

\begin{center}
\small
\begin{tabularx}{\linewidth}{>{\raggedright\arraybackslash}p{0.28\linewidth}YYY}
\toprule
Geom & position [m] & semi-axis [m] & fluidcoef \\
\midrule
\code{fluid\_center\_enclosure} & $(-0.005,\,0,\,0.020)$ & $(0.205,\,0.115,\,0.130)$ & $(0.620,\,0.175,\,1.800,\,0,\,0)$ \\
\code{fluid\_port\_lower\_body} & $(0,\,0.208,\,0.020)$ & $(0.370,\,0.080,\,0.120)$ & $(0.520,\,0.150,\,1.300,\,0,\,0)$ \\
\code{fluid\_starboard\_lower\_body} & $(0,\,-0.208,\,0.020)$ & $(0.370,\,0.080,\,0.120)$ & $(0.520,\,0.150,\,1.300,\,0,\,0)$ \\
\bottomrule
\end{tabularx}
\end{center}

\begin{figure}[H]
  \centering
  \includegraphics[width=0.88\linewidth]{figures_v30/uuv_v30_ellipsoid_layout.png}
  \caption{current scene의 ellipsoid fluid proxy, distributed body component, buoyancy point, CoM, CoB, thruster 방향을 한 그림에 겹쳐 놓은 3D layout.}
\end{figure}

이 그림에서 보듯 current v3.0은 \textbf{fluid drag center, distributed mass component, distributed buoyancy point} 를 함께 쓴다. 따라서 현재 모델은 단순히 ``scene에 ellipsoid geom 세 개 넣었다'' 가 아니라,

\begin{enumerate}[leftmargin=2em]
  \item scene의 ellipsoid geom이 built-in fluid drag를 만들고,
  \item profile의 \code{body\_components} 가 CoM / inertia를 만들고,
  \item profile의 \code{buoyancy\_points} 와 \code{cob\_*\_offset} 이 restoring torque를 만든다.
\end{enumerate}

\subsection{rigid-body / buoyancy / ellipsoid 모델을 식으로 쓰면}

기본 부력은 다음처럼 쓸 수 있다.

\begin{equation}
F_B = \rho g V_{\mathrm{sub}}
\end{equation}

여기서 중요한 것은 current가 \textbf{한 점 부력}이 아니라 \textbf{여러 buoyancy point의 분산 부력}을 쓴다는 점이다. 따라서 전체 부력 토크는 각 점의 합으로 계산된다.

\begin{equation}
\tau_B = \sum_i (r_{B_i} - r_{\mathrm{CoM}}) \times F_{B_i}
\end{equation}

body component를 ellipsoid로 보면 각 component의 자기 관성은 다음과 같이 계산된다.

\begin{align}
I_x &= \frac{m}{5}(b^2 + c^2), \\
I_y &= \frac{m}{5}(a^2 + c^2), \\
I_z &= \frac{m}{5}(a^2 + b^2)
\end{align}

current v3.0에서는 이 합성 관성에 \code{body\_inertia\_scale\_xyz = 1e-6} 을 곱해 runtime inertia를 매우 작게 만든다. 이것이 \textbf{response를 가볍게 만드는 동시에, 강한 fluid drag와 CoB torque scale이 응답을 다시 잡는 구조}다.

또한 built-in ellipsoid fluid는 proxy geom에 지정된 \code{fluidcoef} 와 semi-axis를 기준으로 MuJoCo 내부에서 drag 계열 힘을 계산한다. 따라서 current는 legacy의 직접 damping 행렬과 달리, \textbf{형상 기반(scene 기반) passive fluid term}을 쓴다.

\section{SITL / 제어 파라미터 정합}

depth hold가 실제로 정상 동작하게 만든 것은 물리만이 아니고, \textbf{ArduSub SITL 파라미터 정합}이었다. 현재 시작 스크립트는 다음 값을 강제로 넣는다.

\begin{center}
\small
\begin{tabularx}{\linewidth}{>{\raggedright\arraybackslash}p{0.26\linewidth}YY}
\toprule
파라미터 & 값 & 목적 \\
\midrule
\code{AHRS\_EKF\_TYPE} & 10 & internal EKF 대신 SIM backend를 active AHRS로 사용 \\
\code{AHRS\_GPS\_USE} & 0 & GPS 기반 높이/자세 개입 최소화 \\
\code{GPS1\_TYPE}, \code{GPS2\_TYPE} & 0 & 가상 GPS probing 비활성화 \\
\code{EK3\_IMU\_MASK} & 1 & background EKF3 lane 단일화 \\
\code{EK3\_PRIMARY} & 0 & primary lane 고정 \\
\code{RNGFND1\_TYPE} & 100 & SITL rangefinder backend 사용 \\
\code{RNGFND1\_MIN/MAX} & 0.05 / 30.0 & 수중 하향 rangefinder 거리 범위 \\
\code{RNGFND1\_ORIENT} & 25 & downward-facing \\
\code{THR\_DZ} & 100 & vertical stick deadzone \\
\code{PILOT\_SPEED\_UP/DN} & 180 / 140 & 수직 응답 빠른 preset \\
\code{PILOT\_ACCEL\_Z} & 220 & 수직 가속 응답 \\
\bottomrule
\end{tabularx}
\end{center}

즉 \textbf{depth hold를 살린 결정적 요인}은 ``유체력 모델이 좋아졌다'' 하나가 아니라,

\begin{enumerate}[leftmargin=2em]
  \item vertical estimator를 SIM backend로 고정하고,
  \item fake GPS를 끄고,
  \item rangefinder를 SITL backend로 다시 살리고,
  \item current scene/profile의 restoring torque와 drag를 맞추는
\end{enumerate}

\noindent 조합이었다.

\section{실측 rosbag 기반 legacy vs current 비교}

\subsection{XY path 비교}

\begin{figure}[H]
  \centering
  \includegraphics[width=0.98\linewidth]{figures_v30/uuv_v30_actual_mode_xy_compare.png}
  \caption{실측 rosbag 기반 mode-path XY overlay. current v3.0은 legacy보다 이동량이 훨씬 작고 path가 짧다.}
\end{figure}

\subsection{3D trajectory 비교}

\begin{figure}[H]
  \centering
  \includegraphics[width=0.98\linewidth]{figures_v30/uuv_v30_actual_mode_3d_compare.png}
  \caption{실측 rosbag 기반 3D trajectory 비교. legacy는 더 큰 누적 이동을 보이고, current는 현재 설정상 매우 보수적인 envelope 안에 머문다.}
\end{figure}

\subsection{정량 비교}

\begin{center}
\small
\begin{tabularx}{\linewidth}{>{\raggedright\arraybackslash}p{0.22\linewidth}YYYY}
\toprule
모드 & legacy path [m] & current path [m] & legacy max pitch [deg] & current max pitch [deg] \\
\midrule
MANUAL & 0.3582 & 0.1704 & 21.5138 & 0.1751 \\
STABILIZE & 0.4789 & 0.1083 & 20.0352 & 0.1135 \\
ALT\_HOLD & 0.5138 & 0.1049 & 19.9986 & 0.0341 \\
POSHOLD & 0.5016 & 0.1045 & 20.6876 & 0.0312 \\
\bottomrule
\end{tabularx}
\end{center}

실측 기준으로 보면 current v3.0의 특징은 명확하다.

\begin{itemize}[leftmargin=2em]
  \item \textbf{pitch disturbance가 거의 사라졌다.} 평균 max pitch가 20.56$^\circ$ 에서 0.0885$^\circ$ 로 줄었다.
  \item \textbf{수평 drift가 크게 줄었다.} 평균 horizontal distance가 0.448\,m 에서 0.0273\,m 로 줄었다.
  \item \textbf{응답 authority도 같이 줄었다.} 평균 path length 자체도 0.463\,m 에서 0.122\,m 로 감소했다.
\end{itemize}

즉 current v3.0은 legacy보다 ``더 잘 움직이는'' 모델이라기보다, \textbf{지금 시점에서는 훨씬 안정적이지만 훨씬 보수적인 모델}이다.

\subsection{해석할 때 주의할 점}

이번 actual mode-path 측정에서 90도 yaw turn leg가 충분히 형성되지는 않았다. 따라서 이 절의 yaw 값은 ``90도를 얼마나 잘 따라갔는가'' 보다는 \textbf{실제 achieved yaw magnitude} 로 읽는 것이 맞다. 즉 이번 actual bag는 \textbf{기체 stability와 drift suppression을 보여 주는 근거}로는 매우 유효하지만, large-amplitude maneuver authority 전체를 대표한다고 보기는 어렵다.

\section{current v3.0 actual step 응답}

\begin{figure}[H]
  \centering
  \includegraphics[width=0.98\linewidth]{figures_v30/actual_current_mavros_step_responses.png}
  \caption{current v3.0에서 실제 rosbag로 다시 취득한 RC step / mode step 응답.}
\end{figure}

\begin{center}
\small
\begin{tabularx}{\linewidth}{>{\raggedright\arraybackslash}p{0.32\linewidth}YY}
\toprule
항목 & current actual 값 & 해석 \\
\midrule
manual forward peak $|u|$ & 0.00194\,m/s & forward RC step에도 surge 속도가 거의 생기지 않음 \\
manual yaw peak $|r|$ & 0.000723\,rad/s & yaw step 응답도 매우 작음 \\
stabilize forward peak $|$pitch$|$ & 0.135$^\circ$ & 전진 step 중 자세 disturbance는 매우 작음 \\
stabilize forward peak $|$roll$|$ & 0.123$^\circ$ & 롤 disturbance도 매우 작음 \\
alt-hold heave step depth delta & 0.00114\,m & depth-hold heave 입력에 대한 depth excursion이 1.14\,mm 수준 \\
\bottomrule
\end{tabularx}
\end{center}

\begin{figure}[H]
  \centering
  \includegraphics[width=0.92\linewidth]{figures_v30/actual_current_mavros_mode_comparison.png}
  \caption{current v3.0의 MANUAL / STABILIZE / ALT\_HOLD 구간 비교. 현재 ROS RC path는 큰 움직임보다 안정성과 억제를 더 강하게 보인다.}
\end{figure}

이 결과는 current v3.0이 현재 상태에서 \textbf{high-authority maneuver tuning} 보다는 \textbf{low-drift closed-loop stability} 쪽으로 정렬돼 있음을 보여 준다.

\section{legacy에서 v3.0으로 바꾸며 가장 크게 작용한 것}

\subsection{1. Depth hold를 살린 것은 estimator 경로 정리였다}

실제로 depth hold 무한 하강을 끊은 것은 \code{AHRS\_EKF\_TYPE=10} 과 rangefinder / GPS 정리였다. 즉 \textbf{vertical state 발산을 끊지 않으면} 그 위에 어떤 유체 모델을 올려도 ArduSub는 계속 잘못된 depth target을 추적한다.

\subsection{2. current의 ``안정성''은 ellipsoid-only fluid와 inertia 축소가 크게 만든다}

legacy는 coefficient-based 6-DOF damping / added mass를 직접 계산했다. 반면 current는

\begin{itemize}[leftmargin=2em]
  \item built-in ellipsoid fluid proxy
  \item \code{body\_inertia\_scale\_xyz = 1e-6}
  \item \code{cob\_torque\_scale = 1.0}
  \item 더 큰 linear/angular drag
\end{itemize}

\noindent 조합으로 자세 disturbance를 강하게 억제한다. 실제 bag에서 pitch가 99.57\% 줄어든 것은 이 변화의 직접적인 결과로 읽을 수 있다.

\subsection{3. 통신 정리는 measurement와 운용 재현성을 높였다}

이번 문서용 측정까지 포함하면, 구조적으로 중요한 변화는 다음과 같다.

\begin{itemize}[leftmargin=2em]
  \item JSON sensor port / servo port 혼선 제거
  \item vertical state frame 정합
  \item QGC proxy 도입
  \item measurement 전용 serial2 MAVLink service 경로 복구
  \item readiness / reset / eeprom 유지 최적화
\end{itemize}

즉 current v3.0의 ``안정성''은 단지 물리식 하나의 효과가 아니라, \textbf{estimator + 통신 + scene/profile + runtime force model} 이 같이 정리된 결과다.

\section{현재 모델을 어떻게 설명하는 것이 가장 정확한가}

이번 보고서를 한 문장으로 줄이면 다음과 같다.

\begin{quote}
legacy는 Python이 coefficient-based 6-DOF hydrodynamics를 직접 계산하던 baseline이고, current v3.0은 MuJoCo built-in ellipsoid fluid를 중심에 두되 distributed mass / buoyancy / CoB torque / thrust shaping / SITL bridge를 Python runtime이 보강하는 하이브리드 모델이다.
\end{quote}

그리고 closed-loop actual bag 기준으로는,

\begin{quote}
current v3.0은 legacy보다 훨씬 안정적이고 drift가 적지만, 현재 측정 경로에서는 maneuver authority가 매우 보수적이다.
\end{quote}

이 두 문장을 같이 말하는 것이 가장 정직하고 기술적으로도 맞는 설명이다.

\section{결론}

이번 v3.0 정리에서 중요한 것은 ``그림이 많아졌다''가 아니다. 실제로는 다음이 정리됐다.

\begin{enumerate}[leftmargin=2em]
  \item 실제 rosbag를 다시 채취해서 legacy/current actual 비교 그림을 새로 만들었다.
  \item 전체 SITL--ArduSub--MuJoCo--ROS package 구조와 MuJoCo--SITL 통신 구조를 별도 블록다이어그램으로 정리했다.
  \item current v3.0 물리엔진이 어디까지 MuJoCo이고 어디까지 Python runtime인지, 그리고 ellipsoid가 scene 어디에서 어떤 역할로 쓰이는지 분리해서 설명했다.
  \item depth hold를 살린 estimator 파라미터와, 안정성을 만든 current ellipsoid-only dynamics를 분리해서 설명했다.
\end{enumerate}

지금 시점에서 current v3.0을 평가하면, \textbf{legacy보다 훨씬 안정적이고 실제 closed-loop drift가 적은 쪽}으로 정리됐다고 말할 수 있다. 다만 same bag 기준 maneuver authority는 매우 보수적이므로, 다음 단계는 \textbf{이 안정성을 유지한 채 thrust authority를 얼마나 회복할 것인가}가 된다.

\appendix
\clearpage

\section{부록 A: 보고서 생성에 사용한 파일}

\begin{itemize}[leftmargin=2em]
  \item \code{document/docsource/generate\_uuv\_v30\_doc\_assets.py}
  \item \code{document/docsource/uuv\_mujoco\_v30\_report\_kr.tex}
  \item \code{document/docsource/measurement\_legacy\_mode\_path\_latest\_v30.json}
  \item \code{document/docsource/measurement\_current\_mode\_path\_latest\_v30.json}
  \item \code{document/docsource/measurement\_summary\_current\_latest\_v30.json}
  \item \code{document/docsource/uuv\_v30\_report\_metrics.json}
\end{itemize}

\section{부록 B: 실제 생성된 주요 그림}

\begin{itemize}[leftmargin=2em]
  \item \code{figures\_v30/uuv\_v30\_system\_block\_diagram.png}
  \item \code{figures\_v30/uuv\_v30\_sitl\_comm\_diagram.png}
  \item \code{figures\_v30/uuv\_v30\_physics\_pipeline.png}
  \item \code{figures\_v30/uuv\_v30\_ellipsoid\_layout.png}
  \item \code{figures\_v30/uuv\_v30\_actual\_mode\_xy\_compare.png}
  \item \code{figures\_v30/uuv\_v30\_actual\_mode\_3d\_compare.png}
  \item \code{figures\_v30/uuv\_v30\_actual\_metric\_compare.png}
  \item \code{figures\_v30/actual\_current\_mavros\_step\_responses.png}
\end{itemize}

\end{document}
